\chapter{An\'alisis de requisitos, dise\~no e implementaci\'on}\label{requisitos}
% tfg-floatbarrier-auto
\makeatletter
\@ifundefined{tfgfloatbarrier}{
  \def\tfgfloatbarrier{1}
  \let\tfgorigfigure\figure
  \let\tfgendorigfigure\endfigure
  \renewenvironment{figure}{\tfgorigfigure}{\tfgendorigfigure\FloatBarrier}
  \let\tfgorigtable\table
  \let\tfgendorigtable\endtable
  \renewenvironment{table}{\tfgorigtable}{\tfgendorigtable\FloatBarrier}
}{}
\makeatother

\FloatBarrier
\section{Introducci\'on}

Este cap\'itulo integra tres perspectivas que en un desarrollo profesional no deben
tratarse de forma aislada: la especificaci\'on de requisitos, el dise\~no t\'ecnico y la
implementaci\'on efectiva del sistema.

El objetivo es demostrar trazabilidad de extremo a extremo: desde la necesidad
identificada en el an\'alisis inicial hasta su materializaci\'on en componentes software,
persistencia, seguridad, pruebas y despliegue.

\FloatBarrier
\section{Ecosistema de desarrollo y operaci\'on}

\subsection{Lenguajes y herramientas utilizadas}

El ecosistema t\'ecnico del proyecto se ha definido para garantizar coherencia entre desarrollo, validaci\'on y despliegue, combinando diversas herramientas especializadas en cada \'area del ciclo de vida del software:

\begin{itemize}
    \item \textbf{Backend:} Se emplea \textbf{Java 21} como lenguaje principal, apoyado en el framework \textbf{Spring Boot} (\url{https://spring.io/projects/spring-boot}) para acelerar la creaci\'on de la API REST. La seguridad se gestiona con \textbf{Spring Security}, el acceso a base de datos mediante \textbf{JPA/Hibernate}, la construcci\'on con \textbf{Maven} (\url{https://maven.apache.org/}) y el control de versiones de la base de datos a trav\'es de \textbf{Flyway} (\url{https://flywaydb.org/}).
    \item \textbf{Frontend:} La interfaz de usuario se desarrolla como una Single Page Application utilizando \textbf{React} (\url{https://react.dev/}) y \textbf{TypeScript} (\url{https://www.typescriptlang.org/}) para proveer tipado est\'atico. El empaquetado y entorno de desarrollo r\'apido lo proporciona \textbf{Vite} (\url{https://vitejs.dev/}), usado como \emph{bundler} y servidor de desarrollo para acelerar el arranque y generar el build de producci\'on, mientras que el enrutamiento y las peticiones HTTP se manejan con \textbf{React Router} y \textbf{Axios}. Para el aseguramiento de la calidad se integra \textbf{Vitest} (\url{https://vitest.dev/}) como entorno de pruebas unitarias y \textbf{Stryker} (\url{https://stryker-mutator.io/}) para pruebas de mutaci\'on, introduciendo cambios sint\'eticos en el c\'odigo para medir la fortaleza real de los tests.
    \item \textbf{Documentaci\'on:} La redacci\'on de memorias y documentos formales se realiza en \textbf{\LaTeX{}} (\url{https://www.latex-project.org/}) por su alta calidad tipogr\'afica, apoyado por documentaci\'on auxiliar r\'apida en \textbf{Markdown}.
    \item \textbf{Ingenier\'ia de requisitos y modelado:} Para la gesti\'on de requisitos y la generaci\'on de modelos formales (ERS y DAS) se ha utilizado \textbf{Proteus} (\url{https://github.com/proteus-University-of-Seville/proteus-core}), una herramienta especializada en trazabilidad e ingenier\'ia del software.
    \item \textbf{Gesti\'on del proyecto:} La planificaci\'on iterativa, el backlog y el control del flujo de trabajo se encuentran alojados y versionados en el propio repositorio del proyecto utilizando herramientas nativas de control de versiones.
\end{itemize}

Esta combinaci\'on asegura trazabilidad t\'ecnica, robustez metodol\'ogica y consistencia en los diferentes entornos de despliegue.

\subsection{Pol\'itica de gesti\'on del c\'odigo fuente}

La gesti\'on del c\'odigo fuente de todo el sistema se encuentra centralizada y versionada en un repositorio alojado en GitHub, accesible p\'ublicamente en la direcci\'on \url{https://github.com/Mancalrod/TFG}. 

La evoluci\'on del c\'odigo se realiza mediante Git con un enfoque de ramas tem\'aticas, integraci\'on de nuevas funcionalidades a trav\'es de \emph{pull requests} y la ejecuci\'on de validaciones autom\'aticas previas a su fusi\'on en la rama principal. Esta pol\'itica reduce el riesgo de regresi\'on, mejora significativamente la trazabilidad de los cambios y permite establecer una relaci\'on directa entre los requisitos iniciales, los commits individuales, la ejecuci\'on de pruebas y las etapas de despliegue.

\subsection{Estrategia de integraci\'on continua y despliegue}

La estrategia CI/CD se apoya fundamentalmente en \emph{workflows} orquestados mediante GitHub Actions, divididos en etapas separadas para el backend, el frontend y una validaci\'on conjunta del sistema. En cada cambio relevante detectado en el repositorio, se ejecutan de forma autom\'atica la compilaci\'on, las pruebas unitarias y de integraci\'on, as\'i como an\'alisis de calidad y herramientas de \emph{linting}. A modo ilustrativo, se muestra un fragmento representativo del archivo \texttt{actions.yml} que define estas tareas:

\begin{verbatim}
name: CI Pipeline
on:
  push:
    branches: [ "main", "dev" ]
jobs:
  build:
    runs-on: ubuntu-latest
    steps:
    - uses: actions/checkout@v4
    - name: Set up JDK 21
      uses: actions/setup-java@v4
      with:
        java-version: '21'
        distribution: 'temurin'
    - name: Build with Maven
      run: mvn clean verify
\end{verbatim}

Posteriormente, se construyen las im\'agenes Docker pertinentes con publicaci\'on automatizada y se produce el despliegue hacia una plataforma de Plataforma como Servicio (PaaS) configurada por entornos. En este proyecto se ha optado por desplegar en la plataforma Render, la cual proporciona monitorizaci\'on y gesti\'on en tiempo real del estado de los servicios desplegados. En la Figura \ref{fig:render-deploy} se muestra el panel de control del despliegue en Render.

\figura{0.80}{\detokenize{../Pantallas de TFG/Captura render.png}}{Panel de control y despliegue en la plataforma Render}{fig:render-deploy}{}

El objetivo operativo de este conjunto de medidas es que cada versi\'on candidata sea plenamente reproducible, auditable y desplegable con un nivel de esfuerzo m\'inimo, logrando as\'i una l\'inea de entrega continua s\'olida y compatible con el calendario de desarrollo del TFG.

\subsection{Uso de asistentes de Inteligencia Artificial}

El proyecto adopta un uso de herramientas de IA generativa
limitado al apoyo en tareas de redacci\'on preliminar,
revisi\'on o aceleraci\'on operativa de c\'odigo base,
manteniendo en todo momento el criterio de ingenier\'ia y la autor\'ia de dise\~no por parte del autor.
El alcance y las salvaguardas aplicadas en este \'ambito se detallan
en el Anexo D de esta memoria.

\FloatBarrier
\section{An\'alisis de requisitos}

\subsection{Fuentes de requisitos}

La base fundamental para la definici\'on de los requisitos del proyecto es el propio Documento ERS (Especificaci\'on de Requisitos del Sistema), el cual constituye de por s\'i el cat\'alogo formal y exhaustivo de requisitos. Este documento fue redactado, modelado estructuralmente y versionado haciendo uso de la herramienta Proteus, lo que garantiza rigor metodol\'ogico en su definici\'on.

A partir de este cat\'alogo central establecido en el ERS, el flujo de desarrollo se ha complementado con el Documento DAS (Documento de An\'alisis del Sistema) para el modelado conceptual, as\'i como con el uso din\'amico de un Product Backlog y un Sprint Backlog durante la ejecuci\'on iterativa. La interacci\'on continua con estas herramientas, junto con la evidencia emp\'irica obtenida en las pruebas funcionales y de integraci\'on, asegura que los requisitos no se mantengan como meras definiciones abstractas. En lugar de una especificaci\'on est\'atica, los requisitos del ERS se priorizan, se integran en incrementos de valor y se validan con el comportamiento del producto en funcionamiento, evitando en todo momento cualquier desconexi\'on entre la especificaci\'on te\'orica y la implementaci\'on real del sistema.

\subsection{Estructura del cat\'alogo}

El cat\'alogo se organiza en cuatro bloques complementarios: requisitos funcionales de informaci\'on y operaci\'on, reglas de negocio que acotan el comportamiento permitido, requisitos no funcionales de calidad del sistema y restricciones e integraciones t\'ecnicas por entorno. Esta estructura permite cubrir tanto qu\'e debe hacer el sistema como bajo qu\'e condiciones debe hacerlo.

En Proteus se distinguen expl\'icitamente requisitos funcionales (por ejemplo,
\texttt{RQF-001}), reglas de negocio (\texttt{REGN-001}) y no funcionales
(\texttt{RQNF-001} a \texttt{RQNF-011}), lo que facilita su seguimiento
durante la implementaci\'on.

\subsection{Requisitos funcionales}

El ERS define un abanico funcional amplio, del que se ha priorizado el n\'ucleo
necesario para el MVP. A nivel de dominio, la funcionalidad se agrupa en identidad y acceso (autenticaci\'on de usuarios y control por rol), estructura acad\'emica (cursos, grupos, actividades y vinculaciones), gesti\'on de entregables (configuraci\'on de fechas, visibilidad, reenv\'ios y estructura esperada), gesti\'on de entregas (subida, historial, descarga y trazabilidad) y evaluaci\'on y feedback (calificaci\'on, comentarios y publicaci\'on de notas). Esta agrupaci\'on garantiza cobertura del ciclo acad\'emico completo.

Como ejemplos representativos se incluyen \textbf{RQF-001} (creaci\'on de usuarios por administrador), \textbf{RQF-014--RQF-017} (operaciones de entrega y versionado) y \textbf{RQF-018--RQF-025} (evaluaci\'on, visibilidad e integraciones). Estos ejemplos reflejan la amplitud funcional sin agotar el inventario completo.

\subsection{Reglas de negocio}

Las reglas de negocio restringen de forma expl\'icita qu\'e acciones son admisibles.
Una regla clave es \textbf{REGN-001}, que establece que solo el profesor
autorizado puede asignar una calificaci\'on.

Esta norma se implementa en validaciones de servicio y controlador,
evitando tanto la elevaci\'on de privilegios como la evaluaci\'on fuera de contexto
(p. ej., docente no vinculado al grupo o actividad).

\subsection{Requisitos no funcionales}

Los requisitos no funcionales aportan condiciones de calidad del producto.
La Tabla~\ref{tab:rqnf-resumen} resume los m\'as relevantes para el alcance actual.

\begin{table}[!htbp]
\centering
\begin{tabular}{|P{2.1cm}|P{3.5cm}|P{6.4cm}|}
\hline
\textbf{C\'odigo} & \textbf{Categor\'ia} & \textbf{Aplicaci\'on en el sistema} \\
\hline
RQNF-001 & Fiabilidad & Mitigar p\'erdida de datos en operaciones cr\'iticas de entrega/evaluaci\'on. \\
\hline
RQNF-002 & Usabilidad & Flujos entendibles para usuarios no t\'ecnicos, con navegaci\'on guiada por rol. \\
\hline
RQNF-004 & Mantenibilidad & Bajo acoplamiento entre capas y separaci\'on estricta de responsabilidades. \\
\hline
RQNF-006 & Eficiencia & Respuesta \`agil en consultas frecuentes de panel, actividad y entregables. \\
\hline
RQNF-008 & Portabilidad & Compatibilidad web en navegadores modernos de uso general. \\
\hline
RQNF-009 & Seguridad & Exposici\'on de datos condicionada por autenticaci\'on y autorizaci\'on. \\
\hline
\end{tabular}
\caption{Requisitos no funcionales representativos}
\label{tab:rqnf-resumen}
\end{table}

\subsection{Priorizaci\'on y alcance}

La priorizaci\'on se realiz\'o con criterio MoSCoW y validaci\'on continua con el backlog.
El foco del MVP se situ\'o en los procesos con mayor impacto docente: gesti\'on de actividades y entregables, flujo completo de entrega y evaluaci\'on, seguridad por rol, y persistencia con trazabilidad de cambios. Esta selecci\'on maximiza valor operativo en el horizonte temporal disponible.

Funcionalidades avanzadas (integraciones extendidas, autenticaci\'on federada completa,
observabilidad m\'as profunda) se mantienen en backlog de evoluci\'on y no comprometen
la viabilidad del alcance de TFG.

\subsection{Trazabilidad requisito-backlog-c\'odigo}

La trazabilidad se gestiona en tres niveles estrechamente interconectados: en primer lugar, un \textbf{nivel de an\'alisis} apoyado por el identificador un\'ivoco establecido en el ERS a trav\'es de Proteus; en segundo lugar, un \textbf{nivel de planificaci\'on} guiado por los PBI (\emph{Product Backlog Items}, elementos de trabajo discretos que representan funcionalidades con valor para el usuario dentro de las metodolog\'ias \`agiles) y su sprint asociado; y por \'ultimo, un \textbf{nivel de ejecuci\'on} que liga estos elementos al m\'odulo de software implementado y a la prueba automatizada que lo valida. Esta triple vista permite justificar cualquier decisi\'on de dise\~no con evidencia t\'ecnica medible.

La Tabla~\ref{tab:trazabilidad-requisitos} muestra ejemplos representativos de este encadenamiento, ilustrando c\'omo un requisito formal del ERS se convierte en un PBI dentro de un sprint de desarrollo y, finalmente, se materializa en c\'odigo y pruebas espec\'ificas.

\begin{table}[!htbp]
\centering
\begin{tabular}{|P{2.2cm}|P{3.2cm}|P{4.0cm}|P{2.6cm}|}
\hline
\textbf{Req.} & \textbf{PBI / Sprint} & \textbf{Implementaci\'on principal} & \textbf{Validaci\'on} \\
\hline
RQF-001 & Usuarios y roles / S1-S2 & AuthController, UsuarioService, seguridad por rol & Tests auth + integraci\'on \\
\hline
RQF-018 & Crear entregable / S3-S4 & EntregableController, EntregableService, Entregable & Unit + API tests \\
\hline
RQF-023 & Subir entrega / S4-S5 & EntregaController, EntregaService, MaterialService & Integraci\'on + E2E \\
\hline
REGN-001 & Evaluaci\'on docente / S5 & Validaciones en servicio de entregas y evaluaci\'on & Pruebas de autorizaci\'on \\
\hline
RQNF-009 & Seguridad global / transversal & SecurityConfig, JwtAuthenticationFilter, RateLimitFilter & Tests de seguridad \\
\hline
\end{tabular}
\caption{Ejemplo de trazabilidad entre requisitos, backlog e implementaci\'on}
\label{tab:trazabilidad-requisitos}
\end{table}

\subsection{Contexto documental de partida (ERS y DAS)}

La especificaci\'on documental de este TFG no se apoya en descripciones ad hoc, sino en dos artefactos de an\'alisis formal mantenidos en Proteus. El \textbf{ERS (Especificaci\'on de Requisitos del Sistema)} define situaci\'on actual, necesidades de negocio, subsistemas y cat\'alogo de requisitos (generales, de informaci\'on, funcionales, no funcionales y reglas de negocio), mientras que el \textbf{DAS (Documento de An\'alisis del Sistema)} define arquitectura l\'ogica, modelo de clases, diagramas de estado, diagramas de secuencia, navegabilidad e inventario de mockups de interfaz. Esta base asegura continuidad documental entre an\'alisis y ejecuci\'on.

La reutilizaci\'on de ambos documentos en esta memoria persigue dos objetivos:
primero, preservar trazabilidad respecto al trabajo original de Ingenier\'ia
de Requisitos; segundo, demostrar que las decisiones de dise\~no e implementaci\'on
adoptadas durante el TFG responden a una base de especificaci\'on expl\'icita.

\subsection{Situaci\'on actual y necesidades de negocio (ERS)}

El ERS parte de un an\'alisis detallado de un escenario acad\'emico real donde los procesos docentes presentaban una alta ineficiencia: los profesores requer\'ian m\'ultiples pasos manuales y plataformas dispersas para la definici\'on de actividades, publicaci\'on de material y, de manera muy acusada, para la recogida, versionado y almacenamiento de las entregas de los alumnos. El rastreo de estas entregas carec\'ia de una trazabilidad uniforme (fechas, reenv\'ios, visibilidad), lo que derivaba en p\'erdidas de informaci\'on y en un soporte t\'ecnico que imposibilitaba la integraci\'on de l\'ogicas de evaluaci\'on m\'as avanzadas.

Ante este diagn\'ostico espec\'ifico, las necesidades de negocio que fundamentan este TFG se orientan a un cambio de paradigma centralizado. El objetivo es proporcionar un entorno \'unico donde la carga administrativa del profesor se reduzca mediante flujos estandarizados (desde la creaci\'on de un \emph{entregable} hasta la introducci\'on de calificaciones), y donde el alumno perciba transparencia absoluta sobre sus fechas l\'imite, el estado de sus versiones y los comentarios (feedback) recibidos. Adem\'as, era imperativo garantizar que la evaluaci\'on solo pudiera ser efectuada bajo reglas de autorizaci\'on y contexto acad\'emico estrictas (vinculaci\'on estudiante-grupo-profesor). 

Estas necesidades condicionan el resto de los cap\'itulos en tres planos tangibles: la definici\'on arquitect\'onica (separando estrictamente los servicios de usuarios de la gesti\'on acad\'emica), la implementaci\'on a trav\'es de m\'odulos enfocados a dominios espec\'ificos, y una bater\'ia de pruebas concebida para validar tanto los flujos de experiencia de usuario como el cumplimiento inquebrantable de las reglas de autorizaci\'on impuestas en el negocio.

\subsection{Descomposici\'on en subsistemas (ERS)}

El ERS descompone el sistema en cinco grandes subsistemas l\'ogicos para organizar sus funcionalidades. La Tabla \ref{tab:subsistemas-ers} detalla esta descomposici\'on, la cual ha sido reutilizada como gu\'ia fundamental de dise\~no modular a lo largo del desarrollo del TFG.

\begin{table}[!htbp]
\centering
\small
\begin{tabular}{|P{2.0cm}|P{3.9cm}|P{5.9cm}|}
\hline
\textbf{C\'odigo} & \textbf{Subsistema} & \textbf{Responsabilidad principal} \\
\hline
SUBS-001 & SIU (Interacci\'on con el Usuario) & Portal web y vistas por rol para alumno, profesor y administrador. \\
\hline
SUBS-002 & SGUS (Usuarios y Seguridad) & Autenticaci\'on, autorizaci\'on por rol y control de acceso contextual. \\
\hline
SUBS-003 & SGA (Gesti\'on Acad\'emica) & Cursos, actividades, entregables, evaluaci\'on y feedback. \\
\hline
SUBS-004 & SPA (Procesamiento de Archivos) & Validaci\'on t\'ecnica de ficheros, organizaci\'on y versionado de entregas. \\
\hline
SUBS-005 & SPD (Persistencia de Datos) & Persistencia relacional y repositorio de archivos con trazabilidad. \\
\hline
\end{tabular}
\normalsize
\caption{Subsistemas definidos en ERS y reutilizados en el dise\~no del TFG}
\label{tab:subsistemas-ers}
\end{table}

\subsection{Cat\'alogo exhaustivo de requisitos del ERS}

Con objeto de incorporar de forma completa el contenido de ERS, se documenta
el inventario de requisitos por tipolog\'ia y su interpretaci\'on en el alcance
real del TFG.

\FloatBarrier
\subsubsection*{Requisitos generales (RQG)}

La Tabla \ref{tab:rqg-inventario} detalla los requisitos generales del sistema, que establecen el marco de operaci\'on b\'asico para la gesti\'on de usuarios, actividades, entregables y la comunicaci\'on.

\begin{table}[!htbp]
\centering
\small
\begin{tabular}{|P{1.4cm}|P{3.2cm}|P{7.2cm}|}
\hline
\textbf{C\'odigo} & \textbf{Nombre} & \textbf{Descripci\'on sint\'etica} \\
\hline
RQG-001 & Gesti\'on de usuarios y roles & Alta de usuarios con rol (administrador, profesor, alumno) y permisos diferenciados por perfil. \\
\hline
RQG-002 & Gesti\'on de actividades y entregables & Modelado de actividades evaluables/no evaluables y entregables con reglas propias. \\
\hline
RQG-003 & Acceso de alumnos a entregables & Acceso de estudiante a entregables visibles y registro temporal de acciones. \\
\hline
RQG-004 & Gesti\'on de versiones y reenv\'ios & Control de versiones de entrega y soporte de reenv\'io cuando el docente lo habilita. \\
\hline
RQG-005 & Configuraci\'on de actividades y entregables & Definici\'on de fechas, formatos y condiciones de publicaci\'on por contexto acad\'emico. \\
\hline
RQG-006 & Comunicaci\'on y feedback & Canal de comunicaci\'on docente-estudiante vinculado a la evaluaci\'on de entregas. \\
\hline
RQG-007 & Seguridad y control de acceso & Restricci\'on de visibilidad/operaci\'on por identidad, rol y pertenencia a curso-grupo. \\
\hline
\end{tabular}
\normalsize
\caption{Inventario de requisitos generales (RQG) en ERS}
\label{tab:rqg-inventario}
\end{table}

\FloatBarrier
\subsubsection*{Requisitos de informaci\'on (RQI)}

La Tabla \ref{tab:rqi-inventario} describe los requisitos de informaci\'on que definen qu\'e datos esenciales deben persistirse en el sistema (actividades, entregables, usuarios, etc.).

\begin{table}[!htbp]
\centering
\small
\begin{tabular}{|P{1.4cm}|P{3.2cm}|P{7.2cm}|}
\hline
\textbf{C\'odigo} & \textbf{Nombre} & \textbf{Descripci\'on sint\'etica} \\
\hline
RQI-001 & Informaci\'on sobre actividades & Persistencia de metadatos de actividades creadas por profesorado. \\
\hline
RQI-002 & Informaci\'on sobre entregables & Persistencia de configuraci\'on de cada entregable y su vinculaci\'on con actividad. \\
\hline
RQI-003 & Informaci\'on de usuarios & Registro de identidad, rol y datos de operaci\'on por tipo de usuario. \\
\hline
RQI-004 & Informaci\'on de cursos & Persistencia de estructura acad\'emica (curso/grupo y relaciones). \\
\hline
RQI-005 & Informaci\'on de feedback & Persistencia de calificaciones, comentarios y estado de publicaci\'on. \\
\hline
\end{tabular}
\normalsize
\caption{Inventario de requisitos de informaci\'on (RQI) en ERS}
\label{tab:rqi-inventario}
\end{table}

\subsubsection*{Requisitos funcionales (RQF)}

La Tabla \ref{tab:rqf-inventario-1} y la Tabla \ref{tab:rqf-inventario-2} recogen, de manera estructurada y en dos bloques para facilitar su lectura, las capacidades operativas que el sistema debe ofrecer al usuario final, desde la administraci\'on de identidades hasta la integraci\'on cloud.

\begin{table}[!htbp]
\centering
\small
\begin{tabular}{|P{1.4cm}|P{3.3cm}|P{7.1cm}|}
\hline
\textbf{C\'odigo} & \textbf{Nombre} & \textbf{Descripci\'on sint\'etica} \\
\hline
RQF-001 & Crear usuarios & Alta de usuarios por administrador con asignaci\'on de rol. \\
\hline
RQF-002 & Iniciar sesi\'on & Autenticaci\'on de usuarios en la plataforma. \\
\hline
RQF-003 & Editar usuarios & Modificaci\'on de datos de usuario seg\'un permisos. \\
\hline
RQF-004 & Eliminar usuarios & Baja de usuarios por perfil administrador. \\
\hline
RQF-005 & Crear actividades & Creaci\'on de actividad acad\'emica por docente. \\
\hline
RQF-006 & Visibilizar actividades & Publicaci\'on/ocultaci\'on de actividades para estudiantes. \\
\hline
RQF-007 & Proporcionar adjuntos & Asociaci\'on de material de apoyo a actividad/entregable. \\
\hline
RQF-008 & Modificar actividades & Edici\'on de configuraci\'on de actividad por docente propietario. \\
\hline
RQF-009 & Eliminar actividades & Eliminaci\'on de actividad seg\'un reglas de propiedad. \\
\hline
RQF-010 & Ver actividades & Consulta de actividades por rol y contexto acad\'emico. \\
\hline
RQF-011 & Crear entregables & Alta de entregables dentro de actividad. \\
\hline
RQF-012 & Modificar entregables & Edici\'on de entregable por docente autorizado. \\
\hline
\end{tabular}
\normalsize
\caption{Inventario de requisitos funcionales (RQF-001 a RQF-012)}
\label{tab:rqf-inventario-1}
\end{table}

\begin{table}[!htbp]
\centering
\small
\begin{tabular}{|P{1.4cm}|P{3.3cm}|P{7.1cm}|}
\hline
\textbf{C\'odigo} & \textbf{Nombre} & \textbf{Descripci\'on sint\'etica} \\
\hline
RQF-013 & Eliminar entregable & Baja de entregable por docente autorizado. \\
\hline
RQF-014 & Realizar entrega & Subida de archivo por estudiante para un entregable habilitado. \\
\hline
RQF-015 & Enviar actividad & Confirmaci\'on de env\'io de actividad completa. \\
\hline
RQF-016 & Versionar entregables & Mantenimiento de historial de versiones por entregable. \\
\hline
RQF-017 & Versionar actividades & Mantenimiento de versiones de resoluci\'on de actividad. \\
\hline
RQF-018 & Dar feedback & Publicaci\'on de comentarios docentes sobre entrega/actividad. \\
\hline
RQF-019 & Realizar calificaci\'on & Asignaci\'on de nota por profesor autorizado. \\
\hline
RQF-020 & Ver entregables & Consulta de entregas y versiones por rol permitido. \\
\hline
RQF-021 & Crear curso & Alta de curso/asignatura por administraci\'on. \\
\hline
RQF-022 & Modificar curso & Edici\'on de datos de curso por administraci\'on. \\
\hline
RQF-023 & Vista de estudiante & Visualizaci\'on por profesor de la experiencia de alumno. \\
\hline
RQF-024 & Visibilizar entregable & Publicaci\'on/ocultaci\'on de entregables concretos. \\
\hline
RQF-025 & Integraci\'on cloud & Uso de almacenamiento en nube seg\'un configuraci\'on de entorno. \\
\hline
\end{tabular}
\normalsize
\caption{Inventario de requisitos funcionales (RQF-013 a RQF-025)}
\label{tab:rqf-inventario-2}
\end{table}

\FloatBarrier
\subsubsection*{Reglas de negocio (REGN)}

La Tabla \ref{tab:regn-inventario} agrupa las restricciones ineludibles que condicionan las operaciones del sistema, definiendo estrictamente qui\'en cuenta con los privilegios para ejecutar cada acci\'on, bajo qu\'e contexto acad\'emico particular y respetando qu\'e condiciones temporales o de visibilidad.

\begin{table}[!htbp]
\centering
\small
\begin{tabular}{|P{1.6cm}|P{3.5cm}|P{6.6cm}|}
\hline
\textbf{C\'odigo} & \textbf{Nombre} & \textbf{Regla de negocio} \\
\hline
REGN-001 & Correcci\'on de entregables & Solo el profesor asignado al entregable puede emitir nota. \\
\hline
REGN-002 & Visibilidad de entregables & El alumno solo accede a sus propias entregas/versiones. \\
\hline
REGN-003 & Correcci\'on de actividades & Solo el profesor asignado a la actividad puede calificar. \\
\hline
REGN-004 & Visibilidad de actividades & El alumno solo visualiza sus propias resoluciones de actividad. \\
\hline
REGN-005 & Fecha l\'imite de entrega & No se admiten env\'ios fuera de plazo salvo configuraci\'on expl\'icita. \\
\hline
REGN-006 & Recursos no visibles & Se deniega interacci\'on del alumno con elementos ocultos. \\
\hline
REGN-007 & Creaci\'on de recursos & Solo profesorado crea actividades y entregables. \\
\hline
REGN-008 & Emisi\'on de feedback & Solo el docente autorizado publica feedback evaluativo. \\
\hline
\end{tabular}
\normalsize
\caption{Inventario de reglas de negocio (REGN) en ERS}
\label{tab:regn-inventario}
\end{table}

\FloatBarrier
\subsubsection*{Requisitos no funcionales (RQNF)}

La Tabla \ref{tab:rqnf-inventario-completo} resume las condiciones de calidad t\'ecnica y operacionales que debe cumplir el sistema para garantizar una ejecuci\'on fiable, mantenible y segura, m\'as all\'a de la estricta funcionalidad visible por el usuario.

\begin{table}[!htbp]
\centering
\small
\begin{tabular}{|P{1.6cm}|P{3.2cm}|P{6.9cm}|}
\hline
\textbf{C\'odigo} & \textbf{Nombre} & \textbf{Condici\'on de calidad} \\
\hline
RQNF-001 & Tolerancia a fallos & Persistencia local temporal ante fallos de conectividad para evitar p\'erdida de datos. \\
\hline
RQNF-002 & Intuitividad & Interfaz comprensible para nuevos usuarios con patrones de uso habituales. \\
\hline
RQNF-003 & Uniformidad & Coherencia visual y de estilo con directrices de marca institucional. \\
\hline
RQNF-004 & Acoplamiento bajo & Separaci\'on de responsabilidades con dependencias estrictamente necesarias. \\
\hline
RQNF-005 & Cohesi\'on alta & Organizaci\'on interna de m\'odulos orientada a mantenimiento evolutivo. \\
\hline
RQNF-006 & Tiempo de carga & Tiempo de transici\'on objetivo inferior a 2 segundos en condiciones \`optimas. \\
\hline
RQNF-007 & Tiempo de b\'usqueda & Respuesta r\'apida en consultas frecuentes y minimizaci\'on de b\'usquedas redundantes. \\
\hline
RQNF-008 & Compatibilidad de navegadores & Soporte para navegadores modernos de uso com\'un. \\
\hline
RQNF-009 & Exposici\'on de informaci\'on & Visualizaci\'on de datos restringida a usuarios autenticados y autorizados. \\
\hline
RQNF-010 & Encriptaci\'on & Protecci\'on criptogr\'afica de informaci\'on personal/sensible. \\
\hline
RQNF-011 & Portabilidad web & Compatibilidad expl\'icita con Chrome, Firefox y Edge en escritorio/m\'ovil. \\
\hline
\end{tabular}
\normalsize
\caption{Inventario de requisitos no funcionales (RQNF) en ERS}
\label{tab:rqnf-inventario-completo}
\end{table}

\subsection{Reutilizaci\'on de diagramas y mockups del DAS}

El DAS aporta una base visual amplia (diagramas t\'ecnicos, secuencias,
navegabilidad y mockups) que se reutiliza en esta memoria para explicar
de forma verificable cada decisi\'on de dise\~no e implementaci\'on.

\begin{table}[!htbp]
\centering
\small
\begin{tabular}{|P{3.5cm}|P{5.5cm}|P{2.8cm}|}
\hline
\textbf{Artefacto DAS} & \textbf{Archivo original en assets} & \textbf{Uso en memoria} \\
\hline
Arquitectura l\'ogica & ArquitecturaTFG.drawio.png & Dise\~no global \\
\hline
Diagrama de clases & Diagrama de clases.png & Modelo de dominio \\
\hline
Diagrama de estado & Diagrama de estado.png & Ciclo de vida \\
\hline
Navegabilidad & Diagrama de navegabilidad (1).png & Flujo de interfaz \\
\hline
Secuencia autenticaci\'on & Autenticacion y Acceso.png & Flujo de acceso \\
\hline
Secuencia crear actividad & Backend Activity Creation-2026-03-18-161459.png & Flujo docente \\
\hline
Secuencia realizar entrega & Realizar Entrega.png & Flujo estudiante \\
\hline
Secuencia evaluaci\'on & Evaluacion y Feedback.png & Flujo de correcci\'on \\
\hline
\end{tabular}
\normalsize
\caption{Diagramas principales reutilizados desde DAS}
\label{tab:das-diagramas-principales}
\end{table}

\figura{0.90}{\detokenize{../assets/ArquitecturaTFG.drawio.png}}{Arquitectura l\'ogica del sistema (DAS)}{fig:das-arquitectura}{}

\figura{0.85}{\detokenize{../assets/Diagrama de clases.png}}{Modelo de clases del sistema (DAS)}{fig:das-clases}{}

\figura{0.55}{\detokenize{../assets/Diagrama de estado.png}}{Diagrama de estado de actividades (DAS)}{fig:das-estado}{}

\figura{0.85}{\detokenize{../assets/Diagrama de navegabilidad (1).png}}{Navegabilidad principal de la interfaz (DAS)}{fig:das-navegabilidad}{}

\figura{0.88}{\detokenize{../assets/Autenticacion y Acceso.png}}{Secuencia de autenticaci\'on y acceso (DAS)}{fig:das-seq-auth}{}

\figura{0.88}{\detokenize{../assets/Realizar Entrega.png}}{Secuencia de realizaci\'on de entrega (DAS)}{fig:das-seq-entrega}{}

\figura{0.75}{\detokenize{../assets/InterfazDeUsuarioTFG-LogueadoComoEstudiante.png}}{Mockup de panel principal de estudiante (DAS)}{fig:das-mock-estudiante}{}

\figura{0.75}{\detokenize{../assets/LogueadoComoProfesor.png}}{Mockup de panel principal de profesor (DAS)}{fig:das-mock-profesor}{}

Adicionalmente, el DAS incluye un inventario de diecis\'eis mockups que cubre
las operaciones esenciales por rol. Para mantener equilibrio entre exhaustividad
y legibilidad, su trazado completo se resume en la Tabla~\ref{tab:das-mockups}.

\begin{table}[!htbp]
\centering
\small
\begin{tabular}{|P{4.9cm}|P{6.8cm}|}
\hline
\textbf{Archivo de mockup} & \textbf{Escenario principal representado} \\
\hline
InterfazDeUsuarioTFG-InicioDeSesi\'on.drawio.png & Inicio de sesi\'on y entrada al sistema. \\
\hline
InterfazDeUsuarioTFG-Registrarse.drawio.png & Registro o alta inicial de usuario. \\
\hline
InterfazDeUsuarioTFG-LogueadoComoEstudiante.png & Panel inicial de estudiante autenticado. \\
\hline
VistaEstudianteDentroDeUna\-Asignatura.png & Vista de asignatura desde rol estudiante. \\
\hline
LogueadoComoEstudiante\-Seleccionando\-Checkbox.png & Selecci\'on de entregables por estudiante. \\
\hline
LogueadoComoEstudianteEn\-Calificaciones.png & Consulta de calificaciones y feedback. \\
\hline
LogueadoComoProfesor.png & Panel inicial de profesor autenticado. \\
\hline
VistaProfesorEnAsignatura.png & Vista de asignatura desde rol profesor. \\
\hline
Creacion de actividad.png & Creaci\'on de actividad acad\'emica. \\
\hline
Edici\'on de actividad.png & Edici\'on de actividad existente. \\
\hline
Realizar calificaciones.png & Proceso de calificaci\'on por docente. \\
\hline
DentroDeCalificacionComo\-Profesor.png & Vista detallada de correcci\'on. \\
\hline
EntregaSinNadaPuesto.png & Estado inicial de entrega sin adjuntos. \\
\hline
EntregaConAlgoPuesto.png & Estado de entrega con archivo adjunto. \\
\hline
PonerComentarios.png & Edici\'on de comentarios de evaluaci\'on. \\
\hline
PaginaEstudiantesProfesores.png & Consulta de estudiantes y trazabilidad docente. \\
\hline
\end{tabular}
\normalsize
\caption{Inventario de mockups de interfaz reutilizados desde DAS}
\label{tab:das-mockups}
\end{table}

\subsection{Lectura t\'ecnica de los diagramas estructurales del DAS}

La reutilizaci\'on de las Figuras~\ref{fig:das-arquitectura},
\ref{fig:das-clases}, \ref{fig:das-estado} y \ref{fig:das-navegabilidad}
no se limita a un uso ilustrativo, sino que establece una l\'inea de dise\~no
que condiciona decisiones de implementaci\'on en backend, frontend y datos.

En particular, la arquitectura l\'ogica permite justificar la separaci\'on entre
interacci\'on con usuario, seguridad, gesti\'on acad\'emica, procesamiento de
archivos y persistencia. Este reparto reduce acoplamiento transversal,
facilita pruebas por componente y mejora mantenibilidad evolutiva.

El modelo de clases se emplea como referencia para mantener coherencia entre
entidades persistentes, DTOs y contratos de API. Del mismo modo, el diagrama
de estado de actividades se refleja en validaciones de servicio que evitan
transiciones no permitidas (por ejemplo, publicar, ocultar o cerrar fuera de
las reglas definidas).

Por \`ultimo, el diagrama de navegabilidad refuerza la consistencia entre
estructura funcional y experiencia de usuario por rol, evitando rutas
incoherentes con las restricciones de autorizaci\'on descritas en ERS.

\begin{table}[!htbp]
\centering
\small
\begin{tabular}{|P{2.5cm}|P{3.6cm}|P{3.8cm}|P{2.4cm}|}
\hline
\textbf{Artefacto DAS} & \textbf{Decisi\'on de dise\~no derivada} & \textbf{Materializaci\'on t\'ecnica} & \textbf{Requisitos ERS reforzados} \\
\hline
Arquitectura l\'ogica & Separaci\'on funcional en subsistemas & Frontend SPA + API REST + servicios por dominio + repositorios & SUBS-001 a SUBS-005, RQNF-004 \\
\hline
Diagrama de clases & Normalizaci\'on del dominio acad\'emico & Entidades \texttt{Usuario}, \texttt{Curso}, \texttt{Actividad}, \texttt{Entregable}, \texttt{Entrega} y \texttt{Feedback} & RQI-001 a RQI-005, RQF-005 a RQF-020 \\
\hline
Diagrama de estado & Control expl\'icito del ciclo de vida & Validaciones de transici\'on y reglas temporales en servicios & REGN-005, REGN-006, RQF-006, RQF-024 \\
\hline
Navegabilidad & Flujos de interfaz por rol y contexto & Rutas protegidas, vistas diferenciadas y control de visibilidad & RQG-001, RQG-007, RQNF-002, RQNF-009 \\
\hline
\end{tabular}
\normalsize
\caption{Interpretaci\'on t\'ecnica de los diagramas estructurales del DAS}
\label{tab:das-lectura-estructural}
\end{table}

\subsection{Flujos operativos del DAS y su encaje en la soluci\'on}

Adem\'as de los diagramas estructurales, el DAS aporta una visi\'on de
comportamiento detallada mediante secuencias y mockups operativos. Esta capa
es especialmente relevante para documentar el recorrido completo de profesor y
estudiante en los procesos cr\'iticos del MVP.

El flujo de creaci\'on de actividad docente se explicita en la
Figura~\ref{fig:das-seq-crear-actividad}. En ella se observa la
coordinaci\'on entre interfaz, validaciones de negocio y persistencia,
alineada con RQF-005, RQF-008 y RQF-011.

\figura{0.88}{\detokenize{../assets/Backend Activity Creation-2026-03-18-161459.png}}{Secuencia de creaci\'on de actividad en backend (DAS)}{fig:das-seq-crear-actividad}{}

El ciclo de evaluaci\'on y retroalimentaci\'on queda recogido en la
Figura~\ref{fig:das-seq-evaluacion}, que sirve de base para justificar la
integraci\'on entre calificaci\'on, comentarios y publicaci\'on de resultados,
en coherencia con RQF-018, RQF-019 y REGN-001/REGN-008.

\figura{0.88}{\detokenize{../assets/Evaluacion y Feedback.png}}{Secuencia de evaluaci\'on y feedback (DAS)}{fig:das-seq-evaluacion}{}

En la dimensi\'on de interfaz, los mockups de trabajo muestran c\'omo se
materializa la navegaci\'on por rol en escenarios reales de uso. Las
Figuras~\ref{fig:das-mock-estudiante-asignatura} y
\ref{fig:das-mock-calificacion-profesor} capturan dos momentos clave del
producto: seguimiento de actividad por el estudiante y correcci\'on detallada
por el profesorado.

\figura{0.78}{\detokenize{../assets/VistaEstudianteDentroDeUnaAsignatura.png}}{Mockup operativo: estudiante dentro de una asignatura (DAS)}{fig:das-mock-estudiante-asignatura}{}

\figura{0.78}{\detokenize{../assets/DentroDeCalificacionComoProfesor.png}}{Mockup operativo: detalle de calificaci\'on como profesor (DAS)}{fig:das-mock-calificacion-profesor}{}

Como apoyo adicional, la Figura~\ref{fig:das-mock-entrega-con-archivo}
ilustra el estado de entrega con contenido adjunto, \`util para justificar
la trazabilidad de versiones y el control de evidencias de env\'io.

\figura{0.70}{\detokenize{../assets/EntregaConAlgoPuesto.png}}{Mockup de entrega con archivo adjunto (DAS)}{fig:das-mock-entrega-con-archivo}{}

\begin{table}[!htbp]
\centering
\small
\begin{tabular}{|P{2.1cm}|P{3.4cm}|P{4.0cm}|P{2.8cm}|}
\hline
\textbf{Actor} & \textbf{Operaci\'on central} & \textbf{Evidencia DAS reutilizada} & \textbf{Requisitos ERS asociados} \\
\hline
Administrador & Alta y administraci\'on de identidades/cursos & Mockups de alta de usuario y creaci\'on de curso & RQF-001 a RQF-004, RQF-021, RQF-022 \\
\hline
Profesor & Dise\~no y publicaci\'on de actividad/entregable & Fig. \ref{fig:das-seq-crear-actividad} y mockup de creaci\'on de actividad & RQF-005 a RQF-013, REGN-007 \\
\hline
Estudiante & Entrega y seguimiento de estado & Fig. \ref{fig:das-seq-entrega}, Fig. \ref{fig:das-mock-entrega-con-archivo} & RQF-014 a RQF-017, REGN-005 \\
\hline
Profesor evaluador & Correcci\'on y feedback & Fig. \ref{fig:das-seq-evaluacion}, Fig. \ref{fig:das-mock-calificacion-profesor} & RQF-018, RQF-019, REGN-001, REGN-008 \\
\hline
Profesor y estudiante & Consulta y trazabilidad de resultados & Mockups de comentarios y vista cruzada docente-estudiante & RQF-020, RQI-005, RQNF-009 \\
\hline
\end{tabular}
\normalsize
\caption{Cobertura operativa por rol con evidencia de DAS y requisitos de ERS}
\label{tab:das-cobertura-rol}
\end{table}

\subsection{Cat\'alogo de casos de uso y matriz de trazabilidad ERS-DAS-implementaci\'on}

El DAS define un cat\'alogo amplio de casos de uso; en esta memoria se han
seleccionado cuatro diagramas representativos (Figuras~\ref{fig:das-cu-1} a
\ref{fig:das-cu-4}) que cubren acceso, gesti\'on acad\'emica, entregas y
evaluaci\'on. Esta selecci\'on se corresponde con los flujos implementados en el
MVP y sirve como base visual para enlazar el modelo operativo con la evidencia
de implementaci\'on y pruebas.

Como representaci\'on visual de dicha cobertura, se incluyen cuatro diagramas
de casos de uso de referencia (Figuras~\ref{fig:das-cu-1} a
\ref{fig:das-cu-4}), para documentar el enlace entre el cat\'alogo funcional y
los flujos reales desarrollados.

\figura{0.82}{\detokenize{../assets/CU-1.drawio.png}}{Caso de uso representativo CU-1 (DAS)}{fig:das-cu-1}{}

\figura{0.82}{\detokenize{../assets/CU-2.drawio.png}}{Caso de uso representativo CU-2 (DAS)}{fig:das-cu-2}{}

\figura{0.82}{\detokenize{../assets/CU-3.drawio.png}}{Caso de uso representativo CU-3 (DAS)}{fig:das-cu-3}{}

\figura{0.82}{\detokenize{../assets/CU-4.drawio.png}}{Caso de uso representativo CU-4 (DAS)}{fig:das-cu-4}{}

\FloatBarrier

Para resumir la cobertura por familias operativas dentro del MVP, la
Tabla~\ref{tab:das-familias-casos-uso} agrupa el alcance funcional, el actor
dominante y el grado de cobertura logrado.

\begin{table}[!htbp]
\centering
\small
\begin{tabular}{|P{3.0cm}|P{4.0cm}|P{2.7cm}|P{2.5cm}|}
\hline
\textbf{Familia operativa} & \textbf{Alcance funcional en DAS/ERS} & \textbf{Actor dominante} & \textbf{Cobertura en MVP} \\
\hline
Acceso y sesi\'on & Inicio de sesi\'on, control de identidad, continuidad de sesi\'on & Todos los roles & Completa \\
\hline
Administraci\'on acad\'emica base & Alta/modificaci\'on de cursos, usuarios y contexto de trabajo & Administrador & Completa \\
\hline
Dise\~no docente de actividades & Creaci\'on, edici\'on, visibilidad y configuraci\'on de entregables & Profesor & Completa \\
\hline
Gesti\'on de entregas & Publicaci\'on de trabajo, versionado y consulta de historial & Estudiante & Completa \\
\hline
Evaluaci\'on y feedback & Calificaci\'on, comentarios y visibilidad de resultados & Mixto & Completa \\
\hline
Supervisi\'on y trazabilidad & Vistas cruzadas profesor-estudiante y revisi\'on del estado de entregas & Profesor & Parcial avanzada \\
\hline
\end{tabular}
\normalsize
\caption{Cobertura por familias de casos de uso del DAS dentro del alcance del MVP}
\label{tab:das-familias-casos-uso}
\end{table}

\FloatBarrier

Para reforzar la coherencia documental, la Tabla~\ref{tab:trazabilidad-ampliada}
integra tres niveles de evidencia: requisito formal en ERS, artefacto de
dise\~no en DAS y materializaci\'on en componentes software y pruebas del
proyecto.

\begin{table}[!htbp]
\centering
\small
\begin{tabular}{|P{2.5cm}|P{3.5cm}|P{3.0cm}|P{2.9cm}|}
\hline
\textbf{Bloque ERS} & \textbf{Artefacto DAS reutilizado} & \textbf{Implementaci\'on principal} & \textbf{Evidencia de verificaci\'on} \\
\hline
RQG-001, RQG-007 & Arquitectura + navegabilidad + mockups de sesi\'on & Configuraci\'on de seguridad, filtros JWT, control de rutas & Tests de autenticaci\'on/autorizaci\'on \\
\hline
RQI-001 a RQI-005 & Diagrama de clases + vistas por rol & Entidades JPA, DTOs y contratos API & Pruebas de persistencia e integraci\'on \\
\hline
RQF-005 a RQF-013 & Secuencia de creaci\'on docente + mockups de actividad & Servicios de actividad/entregable y controladores asociados & Unit tests y API tests de gesti\'on docente \\
\hline
RQF-014 a RQF-017 & Secuencia de entrega + estados de mockup & Servicio de entregas, validaciones y versionado & Integraci\'on de entregas y E2E de alumno \\
\hline
RQF-018, RQF-019 & Secuencia de evaluaci\'on + detalle de calificaci\'on & Flujo de feedback/calificaci\'on en backend y frontend & Pruebas funcionales de evaluaci\'on \\
\hline
REGN-001 a REGN-008 & Casos de uso por rol + navegabilidad & Validaciones contextuales en servicios y filtros & Pruebas de reglas de negocio \\
\hline
RQNF-002, RQNF-003 & Inventario de mockups y patrones de interfaz & Dise\~no de pantallas y mensajes de estado & E2E sobre flujos completos \\
\hline
RQNF-004, RQNF-005, RQNF-009 & Arquitectura de subsistemas + separaci\'on de capas & Estructura modular backend/frontend + control de acceso & CI/CD, test suites y auditor\'ia t\'ecnica \\
\hline
\end{tabular}
\normalsize
\caption{Trazabilidad ampliada entre ERS, DAS e implementaci\'on efectiva}
\label{tab:trazabilidad-ampliada}
\end{table}

\FloatBarrier

\subsection{Cobertura documental alcanzada}

Con la incorporaci\'on anterior, este cap\'itulo integra de forma expl\'icita los cinco bloques completos del cat\'alogo ERS (RQG, RQI, RQF, REGN y RQNF), la descomposici\'on en subsistemas usada durante dise\~no e implementaci\'on, los diagramas estructurales y de comportamiento del DAS, el cat\'alogo operativo de casos de uso y familias funcionales del sistema, el conjunto de mockups de interfaz como evidencia de dise\~no funcional y una matriz ampliada de trazabilidad entre especificaci\'on, dise\~no, implementaci\'on y verificaci\'on. Esta integraci\'on evita vac\'ios entre an\'alisis y ejecuci\'on.

De esta forma, la memoria no solo resume resultados: tambi\'en preserva la
trazabilidad documental completa entre especificaci\'on inicial, dise\~no t\'ecnico
y comportamiento implementado.

\FloatBarrier
\section{Dise\~no del sistema}

\subsection{Arquitectura de alto nivel}

La arquitectura adoptada es cliente-servidor desacoplada, con un \textbf{frontend SPA} en React + TypeScript para presentaci\'on y experiencia de usuario, un \textbf{backend REST} en Spring Boot para l\'ogica de negocio y seguridad, y \textbf{persistencia relacional} con JPA/Hibernate y migraciones Flyway. Esta separaci\'on delimita responsabilidades y facilita evoluci\'on modular.

En backend se aplica dise\~no por capas
(controller, service, repository y entity),
que mejora mantenibilidad, testabilidad y capacidad de evoluci\'on.

\subsection{Aportaci\'on arquitect\'onica y tecnol\'ogica}

Desde la perspectiva t\'ecnica, la contribuci\'on principal se concreta en una separaci\'on frontend-backend mediante API REST, una arquitectura por capas en backend que a\'isla responsabilidades, el uso de DTOs para desacoplar contratos de API del modelo de persistencia y una evoluci\'on de esquema controlada por migraciones versionadas. Este dise\~no reduce el acoplamiento transversal, facilita el mantenimiento evolutivo a medio plazo y habilita la incorporaci\'on incremental de nuevas integraciones sin ruptura del n\'ucleo transaccional.

\subsection{Vista de despliegue}

La vista de despliegue se organiza en bloques principales que cubren el \textbf{cliente web} (navegador del usuario final), el \textbf{frontend} como SPA est\'atica servida desde plataforma PaaS (Render), el \textbf{backend} como API REST desplegada en contenedor Docker en Render, los \textbf{datos} en PostgreSQL gestionada accesible solo desde backend, y los \textbf{servicios externos} de almacenamiento cloud, OneDrive y correo transaccional por API HTTP. Esta disposici\'on mantiene una frontera de seguridad clara entre capas.

La separaci\'on de componentes permite escalar y evolucionar cada capa
sin romper el resto del sistema,
manteniendo adem\'as una frontera de seguridad clara entre cliente,
l\'ogica de negocio y persistencia.

\begin{table}[!htbp]
\centering
\begin{tabular}{|P{2.9cm}|P{3.1cm}|P{5.6cm}|}
\hline
\textbf{Nodo} & \textbf{Tecnolog\'ia} & \textbf{Responsabilidad principal} \\
\hline
Cliente & Navegador moderno & Interacci\'on de usuario y consumo de API sobre HTTPS. \\
\hline
Frontend & React + Vite en Render & Renderizado de interfaz, estado de sesi\'on y experiencia por rol. \\
\hline
Backend & Spring Boot (Java 21) en Render & Exposici\'on de endpoints, reglas de negocio, seguridad y orquestaci\'on de persistencia. \\
\hline
Persistencia & Render PostgreSQL & Almacenamiento transaccional y trazabilidad de datos acad\'emicos. \\
\hline
Correo transaccional & Brevo API HTTP & Recuperaci\'on de contrase\~na y notificaciones sin depender de puertos SMTP salientes. \\
\hline
\end{tabular}
\caption{Vista de despliegue simplificada del sistema}
\label{tab:vista-despliegue}
\end{table}

\subsection{Dise\~no de dominio y datos}

El dominio se modela alrededor de la relaci\'on docente-estudiante-actividad, con \texttt{Usuario} como identidad base, \texttt{Profesor} y \texttt{Estudiante} como especializaciones operativas, \texttt{Curso}, \texttt{Grupo} y \texttt{Actividad} para el contexto acad\'emico, \texttt{Entregable} y \texttt{Entrega} para el ciclo de entrega, y \texttt{Feedback}, \texttt{Material} y entidades de soporte (tokens, notificaciones). Esta estructura facilita trazabilidad y coherencia de reglas de negocio.

La persistencia evoluciona de forma versionada mediante migraciones,
lo que permite incorporar nuevas capacidades sin perder consistencia hist\'orica.

\subsection{Dise\~no de API}

La API sigue orientaci\'on REST y separa contratos mediante DTOs. Se definen m\'odulos de endpoint coherentes por contexto de negocio, incluyendo \texttt{/api/auth} para login, refresh y recuperaci\'on de contrase\~na, \texttt{/api/actividades} para gesti\'on de actividades y consulta por usuario, \texttt{/api/entregables} para administraci\'on de entregables, y \texttt{/api/entregas} para subida, revisi\'on, evaluaci\'on y descarga de entregas. Esta segmentaci\'on facilita evoluci\'on de contratos por dominio.

El uso de DTOs evita exponer directamente entidades de persistencia,
reduce acoplamiento y simplifica evoluciones de contrato.

\subsection{Dise\~no de seguridad}

El dise\~no de seguridad incorpora mecanismos complementarios: autenticaci\'on basada en JWT (token de acceso y refresh token), filtros de validaci\'on en la cadena de seguridad, autorizaci\'on por rol y contexto acad\'emico, cifrado de contrase\~nas con BCrypt y limitaci\'on de tasa en operaciones sensibles de autenticaci\'on. Esta combinaci\'on reduce superficie de ataque sin penalizar la experiencia de uso.

Se configura CORS de forma expl\'icita y se diferencia entre rutas p\'ublicas y protegidas,
minimizando superficie de ataque sin comprometer usabilidad.

\subsection{Dise\~no del frontend}

El frontend se organiza en m\'odulos \texttt{pages}, \texttt{components}, \texttt{context} y \texttt{services} con objetivos de experiencia coherente para profesor y estudiante, protecci\'on de rutas seg\'un sesi\'on y rol, gesti\'on centralizada de autenticaci\'on y manejo uniforme de errores y estados de carga. Esta organizaci\'on evita duplicidades y mejora mantenibilidad.

La capa de red encapsula llamadas HTTP y gesti\'on de tokens,
incluyendo renovaci\'on autom\'atica de sesi\'on al expirar el token de acceso.

\FloatBarrier
\section{Implementaci\'on}

\subsection{Implementaci\'on backend}

La implementaci\'on backend se estructura en el paquete principal de la aplicaci\'on, donde destacan \textbf{controladores} como entrada HTTP y validaci\'on de contratos, \textbf{servicios} para l\'ogica de negocio y reglas del dominio, \textbf{repositorios} con acceso a datos mediante Spring Data JPA y un bloque de \textbf{seguridad} con configuraci\'on central y filtros especializados. Esta separaci\'on refuerza la testabilidad y el control de responsabilidades.

Un punto clave es \texttt{EntregaService}, que centraliza validaciones de entrega,
control de versiones, restricciones temporales y operativa de evaluaci\'on,
con apoyo de servicios auxiliares para materiales y notificaciones.

\subsection{Implementaci\'on de flujos cr\'iticos}

Los flujos con mayor impacto funcional son la \textbf{autenticaci\'on y sesi\'on} (login, emisi\'on de tokens, refresco y cierre seguro), la \textbf{publicaci\'on de entregables} (creaci\'on por docente y configuraci\'on de fechas y condiciones), la \textbf{entrega de material} (subida de ficheros, verificaciones, versionado y almacenamiento de metadatos) y la \textbf{evaluaci\'on} (correcci\'on por profesor autorizado, feedback y visibilidad de nota). Esta selecci\'on concentra los riesgos funcionales del sistema.

Cada flujo incorpora validaciones de entrada, control de errores y verificaci\'on de permisos,
asegurando consistencia en escenarios concurrentes y previniendo acciones fuera de alcance.

\subsection{Implementaci\'on frontend}

La implementaci\'on del cliente cubre pantallas y operaciones necesarias para el MVP, incluyendo acceso y gesti\'on de sesi\'on, paneles de actividades y entregables, formularios de entrega y seguimiento, y vista de evaluaciones con retroalimentaci\'on. Esta cobertura garantiza recorrido completo por rol.

Se han introducido medidas pr\'acticas para mejorar usabilidad:
mensajes de validaci\'on temprana, estados visuales de proceso,
y comportamiento resiliente ante errores de red.

\subsection{Persistencia, migraciones e integraciones}

Se utiliza PostgreSQL en escenarios de despliegue y un perfil de desarrollo
local para iteraci\'on r\'apida. Flyway gestiona la evoluci\'on del esquema, incluyendo ampliaciones como atributos para validaci\'on estructural de entregas ZIP, campos para integraci\'on cloud (OneDrive/Cloudinary), tablas y preferencias de notificaci\'on, y tokens de recuperaci\'on de contrase\~na. Esta evoluci\'on versionada protege la consistencia hist\'orica de datos.

Las integraciones externas se dise\~nan por configuraci\'on,
de modo que el sistema puede operar en local sin dependencia obligatoria
de servicios de terceros.
En producci\'on el env\'io de correos se realiza mediante Brevo a trav\'es
de peticiones HTTPS, evitando depender del protocolo SMTP, que puede estar
restringido en plataformas PaaS. Este canal se utiliza para recuperaci\'on
de contrase\~na y notificaciones transaccionales.

\subsection{Verificaci\'on de implementaci\'on}

La verificaci\'on se realiza en capas (unitaria, integraci\'on, E2E y mutaci\'on),
apoyada por pipelines CI/CD y controles de seguridad en preproducci\'on.
Este enfoque proporciona evidencia objetiva de robustez y reduce el riesgo de regresiones.

\FloatBarrier
\section{Discusi\'on t\'ecnica}

\subsection{Decisiones y compromisos}

Las decisiones principales del dise\~no implican compensaciones expl\'icitas. No se busc\'o seleccionar la alternativa m\'as avanzada en cada eje, sino construir una soluci\'on coherente con el alcance del MVP, el tiempo disponible, los requisitos del ERS y la necesidad de mantener trazabilidad entre especificaci\'on, implementaci\'on y pruebas. Por ello, cada decisi\'on se evalu\'o atendiendo a cuatro criterios: valor funcional para el flujo de entregables, mantenibilidad, riesgo de integraci\'on y coste operativo en un entorno acad\'emico.

\subsubsection*{Arquitectura SPA + API}

La separaci\'on entre frontend SPA y backend REST se adopt\'o porque el dominio requiere experiencias diferenciadas por rol y una API reutilizable para operaciones de entrega, evaluaci\'on y administraci\'on. Esta elecci\'on mejora la mantenibilidad y permite evolucionar la interfaz sin modificar el n\'ucleo de negocio. El compromiso asumido es un mayor esfuerzo de orquestaci\'on: gesti\'on de CORS, contratos HTTP, renovaci\'on de sesi\'on y pruebas de integraci\'on entre cliente y servidor. Se consider\'o aceptable porque el proyecto ya incorporaba pipelines, pruebas automatizadas y una estructura de servicios capaz de absorber esa complejidad.

\subsubsection*{Seguridad y gesti\'on de sesi\'on}

El uso de JWT con refresh token se eligi\'o por su ajuste natural a una aplicaci\'on desacoplada y por reducir estado de sesi\'on en servidor. Frente a una sesi\'on server-side tradicional, esta opci\'on facilita despliegues donde el backend puede escalar o reiniciarse sin depender de memoria local de sesi\'on. La contrapartida es que obliga a controlar expiraci\'on, renovaci\'on, cierre de sesi\'on y protecci\'on de rutas con mayor disciplina. Para mitigarlo, se centraliz\'o la validaci\'on en filtros de seguridad y se reforz\'o la autorizaci\'on por rol y contexto acad\'emico.

\subsubsection*{Persistencia, migraciones y almacenamiento}

La persistencia relacional se mantuvo como base del sistema porque el dominio contiene relaciones fuertes entre usuarios, cursos, grupos, actividades, entregables, entregas y feedback. Esta decisi\'on favorece integridad referencial y consultas trazables. Flyway se adopt\'o para que la evoluci\'on del esquema quedara versionada, especialmente al introducir campos de validaci\'on ZIP, preferencias de notificaci\'on e integraciones cloud. Como compromiso, cada cambio de datos exige disciplina de migraci\'on, pero a cambio se reduce el riesgo de divergencias entre entornos.

En almacenamiento de archivos se separaron dos necesidades. Para persistencia operativa de producci\'on se prioriz\'o un proveedor cloud configurable, evitando depender del sistema de archivos ef\'imero de la plataforma PaaS. Para contextos acad\'emicos donde el profesorado ya trabaja con Microsoft 365, se incorpor\'o OneDrive como integraci\'on opcional, no como dependencia obligatoria del n\'ucleo. Esta separaci\'on permite que el sistema funcione sin servicios externos durante desarrollo y que, en producci\'on, pueda activar capacidades institucionales cuando el entorno lo permita.

\subsubsection*{API de OneDrive frente a SharePoint}

La integraci\'on con OneDrive se resolvi\'o mediante Microsoft Graph sobre la unidad del usuario autenticado. Esta decisi\'on responde al problema operativo concreto del proyecto: el profesorado necesita conectar su propia cuenta, explorar carpetas, seleccionar rutas de destino por actividad o por entregable y recibir entregas con una estructura controlada. El acceso directo mediante OAuth2 y permisos delegados permite trabajar sobre la unidad del usuario con operaciones de subida, descarga, listado de carpetas y renovaci\'on de token, manteniendo la integraci\'on acotada al flujo de entregables.

SharePoint se valor\'o como alternativa para escenarios corporativos, pero su orientaci\'on principal es la gesti\'on de sitios, bibliotecas documentales y permisos compartidos a nivel de organizaci\'on. Ese enfoque encaja bien cuando se quiere centralizar documentaci\'on institucional, pero introduce dependencias adicionales de tenant, aprovisionamiento de sitios, bibliotecas y pol\'iticas administradas. Para este TFG, esa complejidad no aportaba una mejora proporcional sobre la necesidad principal: seleccionar carpetas del OneDrive del usuario y organizar entregas de forma flexible. Por ello se descart\'o como v\'ia principal y se prioriz\'o la API de OneDrive mediante Microsoft Graph, dejando SharePoint como posible extensi\'on futura si la implantaci\'on institucional exigiera repositorios compartidos.

\subsubsection*{Despliegue y servicios externos}

El despliegue en PaaS con CI/CD automatizada se adopt\'o para reducir carga operativa y mantener reproducibilidad. Del mismo modo, el correo transaccional mediante Brevo se seleccion\'o por operar sobre API HTTPS, evitando restricciones frecuentes de SMTP saliente en plataformas gestionadas. Estas elecciones no maximizan el control de infraestructura, como ocurrir\'ia con un VPS autogestionado, pero reducen riesgos de administraci\'on y concentran el esfuerzo del TFG en el producto, las pruebas y la trazabilidad funcional.

\subsection{Limitaciones y trabajo futuro}

Aunque el n\'ucleo funcional est\'a cubierto, las decisiones anteriores dejan una hoja de ruta t\'ecnica clara. La autenticaci\'on local con JWT no sustituye a un SSO institucional completo, por lo que una evoluci\'on natural ser\'ia incorporar OAuth2/OIDC, SAML o LDAP como mecanismos de identidad federada. La integraci\'on con OneDrive depende de la configuraci\'on de Azure y de permisos concedidos por el usuario, de modo que un despliegue institucional deber\'ia revisar pol\'iticas de consentimiento, auditor\'ia y retenci\'on de datos. Adem\'as, el sistema valida estructura, extensi\'on y contenido esperado de entregas, pero no incorpora todav\'ia an\'alisis antivirus exhaustivo ni observabilidad avanzada con m\'etricas de uso, tiempos de correcci\'on o alertas operativas.

Estas limitaciones no condicionan la validez del TFG, sino que delimitan una continuidad razonable: reforzar identidad, auditor\'ia, seguridad documental, m\'etricas de producci\'on y requisitos \emph{Should/Could} pendientes sin reescribir el n\'ucleo transaccional.

\FloatBarrier
\section{Conclusiones}

El an\'alisis realizado demuestra coherencia entre lo especificado y lo implementado.
Los requisitos funcionales priorizados se traducen en flujos operativos reales,
los no funcionales se abordan mediante decisiones t\'ecnicas concretas,
y el dise\~no modular facilita mantenimiento, verificaci\'on y evoluci\'on.

Por tanto, el resultado no es solo una implementaci\'on funcional,
sino una soluci\'on con fundamento de ingenier\'ia de software,
alineada con los objetivos acad\'emicos y profesionales del TFG.
